A decoder trained on sequences of length eight fails on sequences of length
sixteen, even when the task is copying. That failure is usually attributed to
positional encodings (PEs): explicit codes place the model at positions it never
saw, and removing them is reported to extrapolate better. We show that this
comparison is confounded by a different decision---where to emit the
end-of-sequence token. With an explicit code the model learns to terminate at
exactly the training length, with EOS probability $1.00$ at that position, so every
longer answer is truncated there; sequence exact match then reports $0.00$ for a
model that computes almost everything correctly. Under a forced output length the
same checkpoints recover $0.91$--$0.99$ token accuracy, and rotary reaches $1.00$.
Removing EOS supervision altogether reverses the ranking: rotary copy is perfect at
$L{=}16$, twice the training length, while \nope{} degrades to $0.43$ token
accuracy. The collapse is therefore not about position ranges---a protocol in which
all positions occur in training does not prevent it---and not about where the code
enters the network: an absolute code injected at $q/k$ behaves identically to one
in the residual stream. It is about addressability of the answer boundary. Exact
match ranks PEs partly by termination policy, and we recommend reporting
forced-length accuracy beside it.
